\section{Related Work}

Long-term memory research has largely evaluated conversational retention,
retrieval, updates, and conflict handling. Mem0 develops a scalable memory
architecture for long-running agent interactions~\cite{chhikara2025mem0}.
LongMemEval evaluates interactive chat memory across extraction, temporal
reasoning, knowledge updates, and abstention~\cite{wu2024longmemeval}.
MemoryAgentBench organizes incremental multi-turn memory around retrieval,
test-time learning, long-range understanding, and selective forgetting~\cite{hu2025memoryagentbench}. MemoryArena extends the focus to interdependent
multi-session agentic tasks~\cite{he2026memoryarena}. EvoArena further models
progressive updates across dynamic environments, including software, and
proposes patch-based evolving memory~\cite{xu2026evoarena}. These works make it
inappropriate to claim that dynamic memory or environment evolution is new.
Agent Workflow Memory separately induces and selectively retrieves reusable
workflows for web-navigation tasks~\cite{wang2024agentworkflowmemory}; it does
not evaluate whether a fresh coding agent produces a correct patch after a
repository evolution.

These evaluations motivate the need to study memory as an intervention in an
agent loop. \bench{} differs in its downstream unit: a precursor engineering
session, a controlled repository evolution, and a fresh agent that must produce
a tested patch. It therefore measures neither a memory-answer score nor a
general preference judgment.

Repository-level benchmarks establish another adjacent line. RepoQA evaluates
long-context code understanding through repository code search~\cite{liu2024repoqa}, while SWE-bench grounds code-agent evaluation in real
GitHub issues~\cite{jimenez2023swebench}. SWE-EVO evaluates long-horizon,
multi-file software evolution~\cite{le2025sweevo}, and ChainSWE evaluates
dependent bug-fix chains in a shared codebase~\cite{jin2026chainswe}. Public repository history is useful
provenance, but it also makes task and solution exposure a methodological
concern; a recent audit of SWE-bench Verified illustrates that risk~\cite{openai2026swebench}. Our artifact treats public history only as a source lead and requires
a separately reviewed private post-snapshot transition before a real repository
can enter the confirmatory corpus.

The distinction is therefore not that \bench{} introduces evolving code or
multi-session memory. It isolates a different intervention: a fresh agent sees
the same evolved repository and ordinary tools in every condition, while only
bounded predecessor evidence changes. The design compares equal-evidence memory
adapters, reports native product behavior separately, and treats stale
predecessor claims as an observable risk rather than assuming that more history
is always useful.
